En un ensayo sintético de una herramienta de corte se registra el desgaste acumulado \(w\), en micrómetros, con \(w(0)=0\) y \(t\) en minutos desde el inicio del ensayo. Se comparan dos propuestas con parámetros constantes no negativos:
\[
\mathrm A:\ w'=a,\qquad \mathrm B:\ w'=bt.
\]
Los tiempos son exactos. Para calibrar se dispone de
\[
\begin{array}{c|cc}t\ (\mathrm{min})&1&2\\\hline
w_{\rm obs}\ (\symup{\mu}\mathrm{m})&1.02&4.02\\
\text{cota de error absoluto}\ (\symup{\mu}\mathrm{m})&0.08&0.10
\end{array}.
\]
Deje fijado por escrito el conjunto de parámetros compatibles de cada propuesta y su predicción para \(w(3)\), antes de usar el dato siguiente.

\emph{Dato reservado:} en \(t=3\) se registra \(8.40\,\symup{\mu}\mathrm{m}\), con error absoluto garantizado no mayor que \(0.10\,\symup{\mu}\mathrm{m}\). Contraste la predicción ya fijada sin recalibrarla con este dato. Decida qué modelo puede mantenerse para el conjunto del ensayo; se admite concluir que ninguno.
